% portada-algebra-lineal-i.tex ───────────────────────────────────────────────
% Portada — Notas de clase, Álgebra Lineal (curso introductorio) — UG
% Punto de partida: la portada original (portada_algebra_lineal.pdf), la cual
%   no tenía fuente editable, así que se reconstruyó su diagrama (dos "órbitas"
%   elípticas y los vectores v1, v2 desde el origen) en TikZ.
% Cambio: se conserva esa ilustración y el contenido (título, autores,
%   profesora, pie), pero se aplica el mismo patrón visual que
%   portada-algebra-lineal-ii.tex / portada_termodinamica.tex (fondo oscuro,
%   retícula de puntos, resplandor ambiental, cifra de volumen gigante
%   translúcida, barra de identificación superior con paleta teal/purple/gold).
% ────────────────────────────────────────────────────────────────────────────

\documentclass[tikz, border=0pt]{standalone}

\usepackage[T1]{fontenc}
\usepackage[utf8]{inputenc}
\usepackage[spanish]{babel}
\usepackage{ebgaramond}
\usepackage[default]{sourcesanspro}
\usepackage{amsmath}
\usepackage{xcolor}
\usetikzlibrary{calc,arrows.meta}

% ── Paleta (misma familia tonal que las demás portadas de "notas de clase") ──
\definecolor{bg}{HTML}{0B0E1A}
\definecolor{panel}{HTML}{141A2E}
\definecolor{grid}{HTML}{FFFFFF}
\definecolor{teal}{HTML}{2DD4BF}
\definecolor{purple}{HTML}{8B7CF6}
\definecolor{gold}{HTML}{F2C14E}
\definecolor{ink}{HTML}{F4F4F5}
\definecolor{muted}{HTML}{8B93A6}

\begin{document}
\begin{tikzpicture}[font=\sffamily]

  \useasboundingbox (0,0) rectangle (210mm,297mm);

  %% ── Fondo ──────────────────────────────────────────────────────────────
  \fill[bg] (0,0) rectangle (210mm,297mm);

  %% ── Resplandor ambiental detrás del diagrama ──────────────────────────
  \foreach \r/\op in {70/0.05, 52/0.06, 34/0.05} {
    \fill[teal, opacity=\op] (98mm,205mm) circle (\r mm);
  }
  \foreach \r/\op in {55/0.04, 36/0.05} {
    \fill[purple, opacity=\op] (128mm,190mm) circle (\r mm);
  }

  %% ── Retícula de puntos ─────────────────────────────────────────────────
  \foreach \x in {8,15,...,203} {
    \foreach \y in {8,15,...,290} {
      \fill[grid, opacity=0.05] (\x mm,\y mm) circle (0.18mm);
    }
  }

  %% ── Cifra de volumen, gigante y translúcida ────────────────────────────
  \node[font=\fontsize{260}{260}\selectfont\rmfamily\bfseries,
        color=panel, anchor=east] at (200mm,88mm) {I};

  %% ── Barra de identificación superior ────────────────────────────────────
  \fill[purple] (14mm,278mm) rectangle (26mm,279mm);
  \fill[teal]   (27mm,278mm) rectangle (35mm,279mm);

  \node[anchor=east, font=\fontsize{8}{8}\selectfont\sffamily,
        color=muted] at (196mm,278.5mm) {NOTAS DE CLASE~\textperiodcentered~VOL.~I};

  %% ── Diagrama: vectores v1, v2 sobre dos "órbitas" elípticas ─────────────
  \coordinate (C) at (105mm,205mm);

  % líneas diagonales tenues que se pierden hacia los bordes (como en la original)
  \foreach \ang/\len in {12/95, 8/70, -170/90, -166/65} {
    \draw[grid, opacity=0.06, line width=0.3pt]
      (C) ++(\ang:12mm) -- ++(\ang:\len mm);
  }

  % ejes tenues
  \draw[grid, opacity=0.14, line width=0.3pt] ($(C)+(-62mm,0)$) -- ($(C)+(62mm,0)$);
  \draw[grid, opacity=0.14, line width=0.3pt] ($(C)+(0,-46mm)$) -- ($(C)+(0,46mm)$);

  % "órbitas": dos circunferencias achatadas (teal) y una elipse rotada (purple)
  \draw[teal, opacity=0.45, line width=0.7pt] (C) ++(0,-2mm) ellipse (30mm and 26mm);
  \draw[teal, opacity=0.30, line width=0.7pt] (C) ++(0,-2mm) ellipse (42mm and 36mm);
  \draw[purple, opacity=0.55, line width=0.8pt, rotate around={-20:(C)}]
    (C) ellipse (48mm and 15mm);

  % vector v1 (gold)
  \draw[gold, line width=1.1pt, -{Latex[length=2.4mm,width=1.8mm]}]
    (C) -- ++(38:52mm)
    node[pos=1, above right=0.5mm, font=\itshape\fontsize{12}{12}\selectfont, color=gold]
    {$v_1$};

  % vector v2 (teal)
  \draw[teal, line width=1.1pt, -{Latex[length=2.4mm,width=1.8mm]}]
    (C) -- ++(207:44mm)
    node[pos=1, below left=0.5mm, font=\itshape\fontsize{12}{12}\selectfont, color=teal]
    {$v_2$};

  \fill[gold] (C) circle (0.9mm);

  %% ── Bloque de título ─────────────────────────────────────────────────
  \fill[gold] (20mm,125mm) rectangle (32mm,126mm);

  \node[anchor=west, font=\fontsize{40}{44}\selectfont\rmfamily\bfseries,
        color=ink] at (19.3mm,109mm) {\'Algebra};
  \node[anchor=west, font=\fontsize{40}{44}\selectfont\rmfamily\bfseries\itshape,
        color=gold] at (19.3mm,90mm) {Lineal};

  \node[anchor=west, font=\fontsize{11}{11}\selectfont\rmfamily\itshape,
        color=muted] at (20mm,80mm) {Notas de clase para un curso introductorio};

  \draw[grid, opacity=0.12, line width=0.3pt] (20mm,70mm) -- (190mm,70mm);

  \node[anchor=west, font=\fontsize{7.5}{7.5}\selectfont, color=muted]
    at (20mm,63.5mm) {ELABORADAS POR};
  \node[anchor=west, font=\fontsize{13}{13}\selectfont\rmfamily,
        color=ink] at (20mm,56mm) {Ricardo Le\'on Mart\'inez};
  \node[anchor=west, font=\fontsize{13}{13}\selectfont\rmfamily,
        color=ink] at (20mm,48.5mm) {Carlos Alberto Rom\'an Rodr\'iguez};

  \draw[grid, opacity=0.09, line width=0.3pt] (20mm,40mm) -- (190mm,40mm);

  \node[anchor=west, font=\fontsize{7.5}{7.5}\selectfont, color=muted]
    at (20mm,33.5mm) {PROFESORA};
  \node[anchor=west, font=\fontsize{13}{13}\selectfont\rmfamily\itshape,
        color=ink] at (20mm,26mm) {Dra.\ Claudia Reynoso Alc\'antara};

  %% ── Pie de página ────────────────────────────────────────────────────
  \fill[teal] (0,0) rectangle (2.4mm,18mm);

  \node[anchor=west, font=\fontsize{7.5}{9}\selectfont, color=muted]
    at (9mm,11mm) {Departamento de Matem\'aticas};
  \node[anchor=west, font=\fontsize{7.5}{9}\selectfont, color=muted]
    at (9mm,7mm) {Divisi\'on de Ciencias Naturales y Exactas~\textperiodcentered~Universidad de Guanajuato};

  \node[anchor=east, font=\fontsize{15}{15}\selectfont\rmfamily,
        color=gold] at (196mm,8.5mm) {2026};

\end{tikzpicture}
\end{document}
